% Generated draft: requires mathematical and source review.
\documentclass[11pt]{article}
\usepackage[margin=25mm]{geometry}
\usepackage{fontspec}
\setmainfont{Latin Modern Roman}
\usepackage{amsmath,amssymb,mathtools,unicode-math}
\setmathfont{Latin Modern Math}
\usepackage{enumitem,needspace,xurl,hyperref}
\setlength{\parindent}{0pt}
\setlength{\parskip}{5pt}
\setlength{\emergencystretch}{3em}
\newcommand{\N}{\mathbb{N}}
\newcommand{\Z}{\mathbb{Z}}
\newcommand{\Q}{\mathbb{Q}}
\newcommand{\R}{\mathbb{R}}
\newcommand{\C}{\mathbb{C}}
\newcommand{\F}{\mathbb{F}}
\newcommand{\A}{\mathbb{A}}
\newcommand{\PP}{\mathbb{P}}
\newcommand{\E}{\mathbb{E}}
\newcommand{\eps}{\varepsilon}
\newcommand{\dd}{\,\mathrm{d}}
\begin{document}
\section*{UnsolvedMath ID: 700909}
\textbf{Source problem code:} catalogue entry
\subsection{Definitions and mathematical statement}
Let \(K\) be a field of characteristic \(2\), and let \(G=\{g_1,\ldots,g_n\}\) be a finite group of odd order \(n\). Let \(X_1,\ldots,X_n\) be algebraically independent over \(K\), put \(F=K(X_1,\ldots,X_n)\), and consider the generic element
\[
x_G=\sum_{j=1}^{n}X_jg_j\in FG.
\]
Its minimal polynomial is the monic polynomial of least positive degree in \(F[\lambda]\) that annihilates \(x_G\). Write it as
\[
\varphi_G(\lambda;X)=\lambda^{m_G}-s_1^G(X)\lambda^{m_G-1}+\cdots+(-1)^{m_G}s_{m_G}^G(X).
\]
The signs are retained to match the source notation, although \(-1=1\) in \(K\). The source calls the last coefficient \(s_{m_G}^G\) the \emph{reduced norm}. It attributes determination of \(G\) by its group determinant to Formanek--Sibley (1990), and determination by this last coefficient to Hoehnke (1991).

Determination by a coefficient is understood with variables indexed by group elements and compared under a bijection of those indexing sets. More explicitly, for another group \(H\) of order \(n\), form \(x_H=\sum_{h\in H}Y_hh\) and its coefficients \(s_i^H\) in the same manner. A bijection \(b:G\to H\) induces substitution \(Y_{b(g_j)}=X_j\).

\textbf{Question (H.-J. Hoehnke, Kourovka Notebook 12.92).} Can the last coefficient in the supplied determination theorem be replaced by some other coefficient or coefficients \(s_i^G\) with \(i<m_G\)? For a proposed single index \(i\), the corresponding determination implication is
\[
\left[\exists\,b:G\to H\text{ bijective such that }
 s_i^G(X)=\left.s_i^H(Y)\right|_{Y_{b(g_j)}=X_j}\right]
\Longrightarrow G\cong H,
\]
for the finite groups of odd order to which that proposed choice applies, with \(i<m_G\) and \(i<m_H\). For a proposed collection of indices, the same bijection must identify every selected coefficient.

\emph{Unresolved scope.} The supplied wording does not specify whether one coefficient or a collection is intended, whether the indices may depend on the group or its minimal-polynomial degree, or whether that degree is retained as comparison data. Thus the displayed implication explains the determination requirement, but the source does not specify a unique quantified choice of replacement coefficients. If \(m_G=1\), there is no positive index \(i<m_G\); this includes the trivial group.
\subsection{Short English statement}
For a finite group of odd order over a field of characteristic \(2\), the source states that the last coefficient of the generic element's minimal polynomial determines the group up to isomorphism. Hoehnke asks whether earlier coefficients can do the same. The intended choice and number of replacement coefficients require clarification.
\subsection{Sources}
\begin{itemize}
\item Local UnsolvedMath registry, permanent ID \(700909\), exact source code \texttt{catalogue entry}, titled \emph{Kourovka Notebook Problem 12.92}. The supplied registry labels the problem \texttt{open}; that label is supplied metadata, not an independently verified assessment.
\item \emph{The Kourovka Notebook (Unsolved Problems in Group Theory)}, Problem 12.92, proposed by H.-J. Hoehnke, 12th issue (1992). The supplied background identifies the consulted edition as the 21st issue edition (2026), arXiv:1401.0300v48. The collection metadata names E. I. Khukhro and V. D. Mazurov as editors.
\item Formanek--Sibley (1990), cited in the supplied statement for determination by the group determinant. No title or further bibliographic details are supplied.
\item Hoehnke (1991), cited in the supplied statement for determination by the coefficient called the reduced norm. No title or further bibliographic details are supplied.
\item G. Frobenius, \emph{Sitzungsber. Preuss. Akad. Wiss. Berlin} (1896), 1343--1382; supplied as a reference for notation.
\item K. W. Johnson, \emph{Math. Proc. Cambridge Phil. Soc.} \textbf{109} (1991), 299--311; supplied as a reference for notation.
\end{itemize}
\end{document}
